\documentclass[12pt]{article}
\usepackage[utf8]{inputenc}
\usepackage[T1]{fontenc}
\usepackage{hyperref}
\title{Prussian Language Resources for K-12 Curriculum Development}
\author{Compiled by Claude Sonnet 5.5 (claude-sonnet-5-5) for Andr\'as Kornai}
\date{}
\begin{document}
\maketitle

\section{Introduction}
Old Prussian, a Baltic language, died out around the early 18th century. A revival movement exists (Prussian Society Tolkemita founded in Germany in 1980; later groups in Lithuania, Poland and Kaliningrad). Reports say several dozen people use revived Prussian, including a few children. The revived language is partly reconstructed, and scholars disagree about its base and methods.

\section{Primary Educational Resources}
\begin{itemize}
\item Pr\={u}siskai: lessons online or in Vilnius, an electronic dictionary, a newspaper, a YouTube channel and recordings: \url{https://prusai.org/about-en.html} and \url{https://www.prusai.org/language-en.html}
\item Wikibooks, Prussian: \url{https://en.wikibooks.org/wiki/Prussian}
\end{itemize}
No school curriculum was found.

\section{Linguistic and Reference Materials}
\begin{itemize}
\item Wikipedia, Old Prussian language: \url{https://en.wikipedia.org/wiki/Old_Prussian_language}
\item Freelang Prussian dictionary (Prussian--English 3,774 words; English--Prussian 1,904; last updated 2014): \url{https://mail.freelang.net/dictionary/prussian.php}
\item Online six-language Prussian dictionary (English, German, Lithuanian, Latvian, Polish, Russian): \url{https://wirdeins.twanksta.org/}
\item Prussian Reconstructions (Vilnius University): \url{https://sitti.vdu.lt/prussian/}
\item Critical discussion of the ``New Prussian'' project: \url{https://sitti.vdu.lt/prussian/princip.htm}
\item Petit, article in Baltistica (2007) on the German influence on the catechism texts: \url{https://www.baltistica.lt/index.php/baltistica/article/download/1179/1101}
\end{itemize}

\section{Religious and Community Resources}
The surviving texts are three Lutheran catechisms printed in K\"onigsberg; the 1561 Enchiridion contains Luther's Small Catechism in full. Source: \url{https://www.ancient-history-sites.com/?p=158386}. Community: Pr\={u}siskai (above).

\section{Digital and Audio Resources}
The Pr\={u}siskai YouTube channel and spoken recordings, and the online dictionaries above.

\section{Curriculum Development Context}
Existing material suits adult learners. For schools, one must choose among competing reconstructions, and the catechism corpus is small and follows German word order. Contact with Pr\={u}siskai would be the first step.

\end{document}